% Part: first-order-logic
% Chapter: model-theory
% Section: partial-iso

\documentclass[../../../include/open-logic-section]{subfiles}

\begin{document}

\olfileid{mod}{bas}{dlo}
\section{بې له سرحدي ټکو ګڼ خطي ترتيبونه}

\begin{defn}
  \emph{بې له سرحدي ټکو ګڼ خطي ترتيب} د هغې !!{language} يو !!a{structure}
  $\Struct{M}$ دے چې يوازينے دوه ځاييز !!{predicate}~$<$ لري او لاندې
  جملې پوره کوي:
  \begin{enumerate}
  \item $\lforall[x][\lnot x < x]$؛
  \item $\lforall[x][\lforall[y][\lforall[z][(x < y \lif (y < z \lif x
    <z ))]]]$؛
  \item $\lforall[x][\lforall[y][(x< y \lor \eq[x][y] \lor y < x)]]$؛
  \item $\lforall[x][\lexists[y][x < y]]$؛
  \item $\lforall[x][\lexists[y][y < x]]$؛
  \item $\lforall[x][\lforall[y][(x < y \lif \lexists[z][(x < z \land
        z < y))]]]$۔
 \end{enumerate}
\end{defn}

\begin{thm}\ollabel{thm:cantorQ}
  هر دوه !!{enumerable}، بې له سرحدي ټکو ګڼ خطي ترتيبونه آيزومورف دي۔
\end{thm}

\begin{proof}
  فرض کړئ $\Struct{M_1}$ او $\Struct{M_2}$ !!{enumerable}، بې له
  سرحدي ټکو ګڼ خطي ترتيبونه دي، ${<_1} = \Assign{<}{M_1}$ او
  ${<_2} = \Assign{<}{M_2}$، او $\PIso{I}$ د دوى تر منځ د ټولو جزوي
  آيزومورفيزمونو سټ دے۔ $\PIso{I}$ تش نۀ دے، ځکه لږ تر لږه
  $\emptyset \in \PIso{I}$۔ ښيو چې $\PIso{I}$ د مخ او شا خاصيت پوره
  کوي۔ بيا $\Struct{M_1} \iso[p] \Struct{M_2}$، او قضيه
  له~\olref[pis]{thm:p-isom1} څخه پايله کېږي۔

  د دې د ښودلو دپاره چې $\PIso{I}$ د مخ خاصيت پوره کوي، فرض کړئ
  $p \in \PIso{I}$ او د~$i = 1$، \dots،~$n$ دپاره
  $p(a_i) = b_i$، او د عموميت له زيان پرته فرض کړئ
  $a_1 <_1 a_2 <_1 \cdots <_1 a_n$۔ يو $a \in \Domain{M_1}$ چې
  راکړل شي، $b \in \Domain{M_2}$ داسې ومومئ:
  \begin{enumerate}
  \item که $a <_1 a_1$ وي، $b \in \Domain{M_2}$ داسې ونيسئ چې
    $b <_2 b_1$؛
  \item که $a_n <_1 a$ وي، $b \in \Domain{M_2}$ داسې ونيسئ چې
    $b_n <_2 b$؛
 \item که د کوم~$i$ دپاره $a_i <_1 a <_1 a_{i+1}$ وي، نو
   $b \in \Domain{M_2}$ داسې ونيسئ چې $b_i <_2 b <_2 b_{i+1}$۔
  \end{enumerate}
  د مطلوب خاصيت لرونکے~$b$ تل موندل کېدے شي، ځکه $\Struct{M_2}$ بې
  له سرحدي ټکو ګڼ خطي ترتيب دے۔ $q = p \cup \{ \langle a, b \rangle
  \}$ داسې تعريف کړئ چې $q \in \PIso{I}$ د~$p$ مطلوبه غځونه وي۔ په
  دې سره د مخ خاصيت ثابت شو۔ د شا خاصيت هم ورته دے۔ نو
  $\Struct{M_1} \iso[p] \Struct{M_2}$؛ او د
  \olref[pis]{thm:p-isom1} له مخې $\Struct{M_1} \iso \Struct{M_2}$۔
\end{proof}

\begin{prob}
  د~\olref[mod][bas][dlo]{thm:cantorQ} ثبوت په دې ښودلو بشپړ کړئ چې
  $\PIso{I}$ د شا خاصيت پوره کوي۔
\end{prob}

\begin{rem}
  $\Struct{S}$، بې له سرحدي ټکو هر !!{enumerable} ګڼ خطي ترتيب،
  ونيسئ۔ بيا (د~\olref{thm:cantorQ} له مخې) $\Struct{S} \iso
  \Struct{Q}$، چې $\Struct{Q} = (\Rat, <)$ هغه !!{enumerable} ګڼ
  خطي ترتيب دے چې د تعريف ساحه ئې د ناطقو عددونو سټ~$\Rat$ دے۔ اوس د
  \olref[thm]{remark:R} !!{structure}~$\Struct{R} = (\Real, <)$ بيا وګورئ۔ مو
  وليدل چې داسې !!a{enumerable} !!{structure}~$\Struct{S}$ شته چې
  $\Struct{R} \elemequiv \Struct{S}$۔ خو $\Struct{S}$، بې له سرحدي
  ټکو يو !!a{enumerable} ګڼ خطي ترتيب دے؛ نو له جوړښت~$\Struct{Q}$
  سره آيزومورف (او ځکه ابتدايي هم‌ارز) !!{structure} دے۔ د ابتدايي هم‌ارزۍ له
  انتقاليت څخه $\Struct{R} \elemequiv \Struct{Q}$ پايله کېږي۔ (دا مو
  په مستقيم ډول هم ښودلے شول: د هماغه مخ او شا دليل په وسيله مو
  $\Struct{R} \iso[p] \Struct{Q}$ ټينګولے شو۔)
\end{rem}
\end{document}
